\textbf{The actual degree-fourteen Bockstein correction.}
Continue the notation $a_7,s_9,n_{11},z_{13}$ of D.2.57, and put
\[
 v_{15}=\operatorname{Sq}^2z_{13}.
 \tag{D.2.58h1}
\]
There is an actual class $\ell_{14}\in H^{14}(L;\FF_2)$ such that
\[
             v_{15}=\operatorname{Sq}^1\ell_{14}.
 \tag{D.2.58h2}
\]
We prove this as an identity of classes on $L$, rather than as a
formal relation in a candidate quotient.

Write $x\in H_6B$ for the label dual to $v_B^2$; thus
$\beta x=d+e$, where $d,e$ are the degree-five labels dual to
$w_B,u_Bv_B$. On $H_{12}T$, let $\ell$ be the coefficient functional
of the ordinary Browder word $\{d,x\}$. Its restriction to $H_{12}E$
annihilates the complete double-suspension kernel in D.2.56:
the five Pontryagin-product directions have no indecomposable
component, and $\{u,D\},\{u,E_6\}$ have ordinary Lie length at
least three, whereas $\{d,x\}$ has length two. Therefore this
functional is the double cohomology suspension of an actual
$\ell_{14}$ on $L$.

One can also see its nonvanishing directly. The class
\[
             r_{12}=\{d,x\}+b\{b,d\}
             \tag{D.2.58h3}
\]
is primitive: the two reduced diagonals are both
$\{b,d\}\otimes b+b\otimes\{b,d\}$ and cancel.
Its image in $B$ is zero, so the actual product splitting puts it
in $H_{12}F$. Every base-generator functional for
$F\to K(\FF_2,4)$ vanishes on it: a composite of homology squares
on its ordinary Lie and product terms cannot produce the bottom
Pontryagin square $a^2$. Primitivity also annihilates all positive
cohomology products. The exact D.2.56 edge consequently places
$r_{12}$ in $H_{12}E$, and $\ell(r_{12})=1$.

The following is an exact identity of functionals on the whole
ambient group $H_{13}T$:
\[
       \operatorname{Sq}^1(\{d,x\}^{*})=(\xi x)^{*}.
       \tag{D.2.58h4}
\]
Indeed the top Bockstein formula gives
$\beta\xi x=\{x,d+e\}$. No ordinary Hall word can contribute a
$\{d,x\}$ coefficient under $\beta$. Lie length is preserved;
in length two, no label has Bockstein $x$, and the only label with
a $d$ component in its Bockstein is $x$ itself. The latter would
require the zero word $\{x,x\}$. Pontryagin products cannot
contribute an indecomposable word. The other restricted generators
in degree thirteen are $\xi f$ and $\xi\{a,b\}$, both with zero
Bockstein. This proves \textup{(D.2.58h4)} using the exact CLM top
formula \cite[Chapter III, Theorem 1.3(3), p.~218]
{I_CohenLadaMay1976}.

Likewise, the even top Nishida formula identifies the restriction
of $(\xi x)^*$ to $H_{13}E$ with the double suspension of
$\operatorname{Sq}^2z_{13}$: the only degree-six input whose
Bockstein has a $d$ component is $x$.
Thus \textup{(D.2.58h4)} proves the desired equality after double
cohomology suspension. It remains to justify injectivity in this
one additional degree.

The two actual bar constructions give that injectivity without
assuming a new Eilenberg--Moore edge range. In total degree fourteen
of the first bar construction, the only possible filtration-two
cycles beyond filtration one are
\[
 u|S+S|u,\qquad u|R+R|u,\qquad D|E_6+E_6|D.
 \tag{D.2.58h5}
\]
They are the leading filtration terms of actual ordered products
in $H_{14}(\Omega L)$ and hence are killed by the next homology
suspension. In total degree fifteen of the second bar construction,
the four possible filtration-two terms have factors of degrees
six and seven. Their differentials are precisely the four
independent ordered products proved in D.2.56. Hence they supply
no additional class. Filtration three starts above both ranges,
so no higher differential changes these conclusions. Therefore
\[
       H_{13}E\longrightarrow H_{15}L
       \quad\hbox{is surjective}.
       \tag{D.2.58h6}
\]
Dualizing proves the required injectivity and establishes
\textup{(D.2.58h2)} on the actual space $L$.

In particular every reduction of an integral homology class in
degree fifteen annihilates $v_{15}$. This is the ordinary
Bockstein pairing consequence of \textup{(D.2.58h2)}; it does not
assert that every mod-two Bockstein cycle has an integral lift.
